% =============================================================================
% Dyskusja i wnioski. Plik edytowany ręcznie -- skrypt analiza_odpowiedzi_naukowa.py go NIE nadpisuje.
% Liczby: definicje.tex; literatura: references.bib.
% =============================================================================
\section{Dyskusja}

\subsection{Główne wyniki}
W ciągu trzech lat rutynowego stosowania programu mikrouczenia w miejscu pracy odsetek poprawnych
odpowiedzi na wielokrotnie prezentowane pytania rósł wraz z liczbą przypomnień zgodnie z funkcją
potęgową, a efekt był odporny na ubywanie obserwacji, specyfikację modelu i klastrowanie
danych. Poprawność malała wraz z odstępem od poprzedniego przypomnienia, wcześniejsze błędy były
w kolejnej próbie korygowane tylko częściowo, a powtarzane błędy w większości odtwarzały ten sam
dystraktor. Poprawność odpowiedzi na nowe pytania nie rosła wraz ze stażem. % SPRAWDŹ

\subsection{Porównanie z wcześniejszymi badaniami}
Przebieg przyrostu odpowiada krzywym uczenia opisanym dla wielu zbiorów danych
edukacyjnych \citep{koedinger2023} i jest zgodny z efektami testowania i rozłożenia w czasie
\citep{roediger2006,cepeda2006}. Randomizowane badania edukacji rozłożonej w czasie wykazały lepsze
utrwalenie wiedzy \citep{kerfoot2007mededuc,kerfoot2009}; przedstawione dane obserwacyjne rozszerzają te
ustalenia na niemedyczne środowisko pracy z długimi odstępami między powtórkami. Ujemny związek między
odstępem a poprawnością jest zgodny z modelami krzywej zapominania dopasowywanymi do logów uczących się
\citep{settles2016,tabibian2019}. Podobnie jak w badaniach, w których wielokrotnie wykorzystywano te
same pytania \citep{donker2022}, poprawy w powtarzanych pytaniach nie da się oddzielić od zapamiętania
konkretnego pytania i jego odpowiedzi; brak efektu stażu w nowych pytaniach sugeruje ograniczony
transfer.

\subsection{Implikacje praktyczne}
Program ponownie prezentował pytania w tym samym odstępie niezależnie od wyniku poprzedniej próby.
Harmonogramy adaptacyjne, które skracają odstęp po błędach i wydłużają go po poprawnych odpowiedziach,
efektywniej wykorzystują historię uczącego się \citep{lindsey2014,settles2016,tabibian2019}. Wybory
tych samych dystraktorów utrzymywały się, mimo że po każdej odpowiedzi wyświetlano wyjaśnienie poprawnej odpowiedzi,
co sugeruje, że ogólne wyjaśnienie nie wystarcza do skorygowania błędnych przekonań; skuteczniejsza może
być informacja zwrotna odnosząca się do konkretnego wybranego dystraktora. Analiza pytań wskazała
kandydatów do zmiany: pytania o niskiej dyskryminacji, pytania, na które niemal wszyscy odpowiadali
poprawnie już w pierwszej próbie (mało miejsca na uczenie się), oraz dystraktory niefunkcjonalne
\citep{tarrant2009,gierl2017}.

\subsection{Ograniczenia}
Badanie ma charakter obserwacyjny i nie obejmuje grupy kontrolnej, nie można więc przypisać przyrostu
wiedzy programowi w sensie przyczynowym; poprawa w powtarzanych pytaniach może częściowo odzwierciedlać
rozpoznawanie pytania, a nie jego zrozumienie. Udział i aktywność różniły się między pracownikami,
a wyniki dla sekwencji z wieloma próbami opisują bardziej aktywnych uczestników, choć analizy
zbalansowanej kohorty i długości sekwencji nie wskazały na obciążenie wynikające z ubywania obserwacji.
Odsetek poprawnych odpowiedzi w pierwszych próbach (\AccFirst\%) ogranicza zakres możliwej poprawy
(efekt sufitowy).
Dystraktory, których nigdy nie wybrano, nie są widoczne w logu, więc funkcjonowanie dystraktorów może
być przeszacowane. Nie dysponowano informacjami o stanowiskach pracowników ani ich wcześniejszej wiedzy,
a dane o wieku i płci były niedostępne dla \DemoUnmatched{} z \DemoN{} pracowników; cech
demograficznych nie uwzględniano w modelach. Znaczniki czasu nie zawierały informacji o strefie
czasowej. Model mieszany zawierał wyłącznie losowe wyrazy wolne; nie modelowano indywidualnego tempa uczenia się pracowników.
Dane pochodzą z jednej organizacji.

\section{Wnioski}
Powtarzanemu przypominaniu w ramach rutynowych szkoleń pracowniczych towarzyszyła poprawa ograniczona
do powtarzanych pytań oraz zapominanie między powtórkami. Dane z logów takich programów
mogą służyć do monitorowania uczenia się, oceny pytań i projektowania adaptacyjnych harmonogramów
powtórek.
